\begin{figure}[H]
\centering
\begin{tikzpicture}[>=Latex,x=1.10cm,y=0.95cm]
    \def\ang{24}
    \coordinate (O) at (0,0);
    \coordinate (Xp) at (\ang:4.75);
    \coordinate (Tp) at (90-\ang:4.75);

    \fill[minkfieldred] (O) -- (Xp) -- ($(Xp)+(Tp)$) -- (Tp) -- cycle;
    \foreach \i in {-4,-3,...,5}{
        \draw[mink grid x] (\i*0.7,-0.8) -- (\i*0.7,5.0);
        \draw[mink grid t] (-4.15,\i*0.7) -- (4.85,\i*0.7);
        \draw[mink grid prime] (\ang:{\i*0.66}) -- ++(90-\ang:4.65);
        \draw[mink grid prime] (90-\ang:{\i*0.66}) -- ++(\ang:4.65);
    }
    \draw[mink axis] (0,-0.85) -- (0,5.15) node[left] {$ct$};
    \draw[mink axis] (-4.2,0) -- (5.0,0) node[below] {$x$};
    \draw[mink axis prime] (90-\ang:-0.8) -- (Tp) node[right] {$ct'$};
    \draw[mink axis prime] (\ang:-0.8) -- (Xp) node[below right] {$x'$};

    \coordinate (A) at ($(\ang:1.0)+(90-\ang:0.70)$);
    \coordinate (B) at ($(A)+(90-\ang:3.35)$);
    \coordinate (C) at (B |- A);

    \fill[minkpurple,opacity=0.10] (A) -- (B) -- (C) -- cycle;
    \draw[minkdarkred,line width=2.2pt,line cap=round] (A) -- (B)
        node[midway,left=5pt] {\contour{white}{$c\Delta\tau$}};
    \draw[minkdarkblue,line width=1.8pt,line cap=round] (B) -- (C)
        node[midway,right=5pt] {\contour{white}{$c\Delta t$}};
    \draw[minkdarkblue,line width=1.8pt,line cap=round] (C) -- (A)
        node[midway,below=4pt] {\contour{white}{$\Delta x$}};
    \draw[minkdarkblue,semithick]
        ($(C)+(-0.20,0)$) -- ($(C)+(-0.20,0.20)$) -- ($(C)+(0,0.20)$);

    \fill[minkdarkred] (A) circle (2.3pt)
        node[below left=2pt] {\contour{white}{tic}};
    \fill[minkdarkred] (B) circle (2.3pt)
        node[above left=2pt] {\contour{white}{toc}};
    \filldraw[fill=white,draw=minkdarkblue,line width=1pt] (C) circle (2.1pt);

    \node[minkdarkred,font=\small,align=left,anchor=west] at (2.95,4.55)
        {reloj fijo en $S'$\\$\Delta x'=0$};
    \draw[mink small arrow,minkdarkred] (2.85,4.18) -- (B);
    \node[font=\small,align=left,anchor=west] at (-3.85,4.55)
        {$\Delta t=\gamma\Delta\tau$\\el intervalo de $S$ es mayor};
\end{tikzpicture}
\caption{Dilatación temporal para un reloj en reposo en $S'$. El lado rojo inclinado une tic y toc y representa el tiempo propio $\Delta\tau$; el triángulo coordenado muestra la separación espacial $\Delta x$ y el intervalo mayor $\Delta t$ medido desde $S$.}
\label{fig:dilatacion_Sp}
\end{figure}